% GENERATED FILE. Edit canonical lessons, then run npm run generate:books.
% canonical-source-hash: fnv1a64:73832008f1146077
% canonical-lessons: LA-C294-identidem, LA-C294-subinde, LA-C294-crebro, LA-C294-plerumque, LA-C294-quotiens

\chapter{Identidem}
\label{ch:la-identidem}
\hlchaptermodality{294}

\begin{chapteropening}
\textbf{By the end of this chapter:} \emph{I can say how often things happened: again and again, immediately after, frequently, mostly and as often as.}

Complete the last lesson of chapter 294: I can say how often things happened: again and again, immediately after, frequently, mostly and as often as.
\end{chapteropening}

\section[identidem]{identidem — again and again, repeatedly}

\label{lesson:LA-C294-identidem}



\begin{quote}
Before the new one: say the Latin for for a little while, then the Latin for for some time.
\end{quote}

\subsection*{You'll want to know: identidem}
\textbf{identidem} — \textquotedblleft{}again and again, repeatedly\textquotedblright{}.

\begin{quote}
Īnfāns identidem flēbat.
\end{quote}

\begin{quote}
\emph{The baby cried again and again.}
\end{quote}

An adverb for an action that keeps coming back, like \emph{saepe} (taught) but closer together.

\begin{cousinweb}
\textbf{identidem} is built on \emph{idem}, \textquotedblleft{}the same\textquotedblright{} (EtymDB): the same, and the same again. Lewis \& Short analyse the middle part in more than one way.
\end{cousinweb}

\begin{grammarlens}[title={What you've built}]
One more word for telling, explaining and describing at B1.
\end{grammarlens}

\subsection*{Your turn}
\begin{itemize}
\raggedright
  \item \emph{Say it:} \emph{identidem}
  \item \emph{Say it:} \emph{identidem}, once more
  \item \emph{Say it:} a sentence of your own with \emph{identidem}
  \item \emph{Recall:} say the Latin for for a little while, then the Latin for for some time, then say \emph{identidem} again
\end{itemize}

\begin{tcolorbox}[breakable,colback=teal!4,colframe=teal!35!black,title={Before you move on}]
What does identidem mean? (\textquotedblleft{}Again and again, repeatedly\textquotedblright{}.) Say it once more.
\end{tcolorbox}
\section[subinde]{subinde — immediately after; now and then}

\label{lesson:LA-C294-subinde}



\begin{quote}
Before the new one: say the Latin for for some time, then the Latin for again and again.
\end{quote}

\subsection*{You'll want to know: subinde}
\textbf{subinde} — \textquotedblleft{}immediately after; now and then\textquotedblright{}.

\begin{quote}
Prīmō rīsit, subinde flēvit.
\end{quote}

\begin{quote}
\emph{At first she laughed, and immediately after she wept.}
\end{quote}

Horace and later prose. In Seneca and Pliny it also means 'from time to time'.

\begin{cousinweb}
\textbf{subinde} joins \emph{sub}, \textquotedblleft{}close after\textquotedblright{}, and \emph{inde}, \textquotedblleft{}from there, then\textquotedblright{} (Lewis \& Short): following closely on.
\end{cousinweb}

\begin{grammarlens}[title={What you've built}]
One more word for telling, explaining and describing at B1.
\end{grammarlens}

\subsection*{Your turn}
\begin{itemize}
\raggedright
  \item \emph{Say it:} \emph{subinde}
  \item \emph{Say it:} \emph{subinde}, once more
  \item \emph{Say it:} a sentence of your own with \emph{subinde}
  \item \emph{Recall:} say the Latin for for some time, then the Latin for again and again, then say \emph{subinde} again
\end{itemize}

\begin{tcolorbox}[breakable,colback=teal!4,colframe=teal!35!black,title={Before you move on}]
What does subinde mean? (\textquotedblleft{}Immediately after; now and then\textquotedblright{}.) Say it once more.
\end{tcolorbox}
\section[crēbrō]{crēbrō — frequently, often}

\label{lesson:LA-C294-crebro}



\begin{quote}
Before the new one: say the Latin for again and again, then the Latin for immediately after.
\end{quote}

\subsection*{You'll want to know: crēbrō}
\textbf{crēbrō} — \textquotedblleft{}frequently, often\textquotedblright{}.

\begin{quote}
Ad vīllam avī crēbrō veniēbāmus.
\end{quote}

\begin{quote}
\emph{We often used to come to our grandfather's villa.}
\end{quote}

The adverb of \emph{crēber} (taught), 'thick, close-packed, frequent'.

\begin{cousinweb}
\textbf{crēbrō} is the ablative of \emph{crēber}, \textquotedblleft{}close-packed, frequent\textquotedblright{}, used as an adverb (EtymDB).
\end{cousinweb}

\begin{grammarlens}[title={What you've built}]
One more word for telling, explaining and describing at B1.
\end{grammarlens}

\subsection*{Your turn}
\begin{itemize}
\raggedright
  \item \emph{Say it:} \emph{crēbrō}
  \item \emph{Say it:} \emph{crēbrō}, once more
  \item \emph{Say it:} a sentence of your own with \emph{crēbrō}
  \item \emph{Recall:} say the Latin for again and again, then the Latin for immediately after, then say \emph{crēbrō} again
\end{itemize}

\begin{tcolorbox}[breakable,colback=teal!4,colframe=teal!35!black,title={Before you move on}]
What does crēbrō mean? (\textquotedblleft{}Frequently, often\textquotedblright{}.) Say it once more.
\end{tcolorbox}
\section[plērumque]{plērumque — mostly, usually}

\label{lesson:LA-C294-plerumque}



\begin{quote}
Before the new one: say the Latin for immediately after, then the Latin for frequently.
\end{quote}

\subsection*{You'll want to know: plērumque}
\textbf{plērumque} — \textquotedblleft{}mostly, usually\textquotedblright{}.

\begin{quote}
Plērumque post merīdiem dormiēbat.
\end{quote}

\begin{quote}
\emph{He usually slept in the afternoon.}
\end{quote}

The neuter of \emph{plērusque}, 'the greater part': for the most part.

\begin{cousinweb}
\textbf{plērumque} is the neuter of \emph{plērusque}, \textquotedblleft{}the greater part, most\textquotedblright{} (Lewis \& Short), used as an adverb: for the most part.
\end{cousinweb}

\begin{grammarlens}[title={What you've built}]
One more word for telling, explaining and describing at B1.
\end{grammarlens}

\subsection*{Your turn}
\begin{itemize}
\raggedright
  \item \emph{Say it:} \emph{plērumque}
  \item \emph{Say it:} \emph{plērumque}, once more
  \item \emph{Say it:} a sentence of your own with \emph{plērumque}
  \item \emph{Recall:} say the Latin for immediately after, then the Latin for frequently, then say \emph{plērumque} again
\end{itemize}

\begin{tcolorbox}[breakable,colback=teal!4,colframe=teal!35!black,title={Before you move on}]
What does plērumque mean? (\textquotedblleft{}Mostly, usually\textquotedblright{}.) Say it once more.
\end{tcolorbox}
\section[quotiēns]{quotiēns — as often as; how often}

\label{lesson:LA-C294-quotiens}



\begin{quote}
Before the new one: say the Latin for frequently, then the Latin for mostly.
\end{quote}

\subsection*{You'll want to know: quotiēns}
\textbf{quotiēns} — \textquotedblleft{}as often as; how often\textquotedblright{}.

\begin{quote}
Quotiēns in urbem veniēbat, amīcōs vīsitābat.
\end{quote}

\begin{quote}
\emph{As often as he came into the city, he visited his friends.}
\end{quote}

Its partner is \emph{totiēns}, 'so often'. English \emph{quotient} comes from it.

\begin{cousinweb}
\textbf{quotiēns} is built on \emph{quot}, \textquotedblleft{}how many\textquotedblright{} (Lewis \& Short): how many times. English \emph{quotient} was formed from \emph{quotiēns} (EtymDB).
\end{cousinweb}

\begin{grammarlens}[title={What you've built}]
One more word for telling, explaining and describing at B1.
\end{grammarlens}

\subsection*{Your turn}
\begin{itemize}
\raggedright
  \item \emph{Say it:} \emph{quotiēns}
  \item \emph{Say it:} \emph{quotiēns}, once more
  \item \emph{Say it:} a sentence of your own with \emph{quotiēns}
  \item \emph{Recall:} say the Latin for frequently, then the Latin for mostly, then say \emph{quotiēns} again
\end{itemize}

\begin{tcolorbox}[breakable,colback=teal!4,colframe=teal!35!black,title={Before you move on}]
What does quotiēns mean? (\textquotedblleft{}As often as; how often\textquotedblright{}.) Say it once more.
\end{tcolorbox}
